\documentclass[11pt,a4paper]{article}

% ---------- Datos del documento ----------
\newcommand{\titulo}{Resolución de una ecuación de primer grado con la balanza algebraica de Polypad}
% Título corto para el encabezado de las páginas
\newcommand{\tituloCorto}{La balanza: $5x+12=67$}
\newcommand{\asignatura}{Didáctica de las Matemáticas}
\newcommand{\autor}{Javier Gómez López}
% Texto explicativo de la tarea (déjalo vacío para que no aparezca)
\newcommand{\enunciado}{%
Resolver la ecuación $5x + 12 = 67$ utilizando la balanza algebraica de
Polypad y explicar cada uno de los pasos realizados, relacionando lo que
ocurre en la balanza con la manipulación algebraica de la ecuación.%
}
\newcommand{\master}{MÁSTER UNIVERSITARIO EN FORMACIÓN DEL PROFESORADO-ESPECIALIDAD EN MATEMÁTICAS -- Universidad Rey Juan Carlos}

% ---------- Paquetes ----------
\usepackage[utf8]{inputenc}
\usepackage[T1]{fontenc}
\usepackage[spanish,es-noshorthands]{babel}
\usepackage{lmodern}
\usepackage[margin=2.5cm,headheight=14pt]{geometry}
\usepackage{xcolor}
\usepackage{titlesec}
\usepackage{fancyhdr}
\usepackage{graphicx}
\usepackage{enumitem}
\usepackage{parskip}
\usepackage{etoolbox}
\usepackage{tabularx}
\usepackage{booktabs}
\usepackage{amsmath,amssymb}
\usepackage{float}
\usepackage[font=small,labelfont={bf,color=gris}]{caption}
\usepackage[most]{tcolorbox}
\usepackage[hidelinks]{hyperref}

% ---------- Estilo ----------
\definecolor{acento}{RGB}{176,0,32}   % rojo URJC
\definecolor{gris}{RGB}{90,90,90}

\titleformat{\section}{\Large\bfseries\color{acento}}{\thesection.}{0.6em}{}
\titleformat{\subsection}{\large\bfseries\color{gris}}{\thesubsection.}{0.6em}{}

\pagestyle{fancy}
\fancyhf{}
\fancyhead[L]{\small\color{gris}\asignatura}
\fancyhead[R]{\small\color{gris}\tituloCorto}
\fancyfoot[C]{\small\thepage}
\renewcommand{\headrule}{{\color{acento}\hrule height 0.6pt}}


% Captura de la balanza
\newcommand{\captura}[2]{%
  \begin{figure}[H]
    \centering
    \includegraphics[width=0.72\linewidth]{#1}
    \caption{#2}
  \end{figure}}

\begin{document}

% ---------- Cabecera ----------
\thispagestyle{empty}
\begin{center}
  {\small\color{gris}\MakeUppercase{\master}}\\[1.2em]
  {\color{acento}\rule{\linewidth}{1.2pt}}\\[0.8em]
  {\Huge\bfseries \titulo}\\[0.6em]
  {\Large\color{gris} \asignatura}\\[0.6em]
  {\color{acento}\rule{\linewidth}{1.2pt}}\\[1em]
  {\large \autor}\\[0.3em]
  {\small\color{gris}\today}
\end{center}
\vspace{1.5em}

% ---------- Enunciado ----------
\ifdefempty{\enunciado}{}{%
  \begin{tcolorbox}[enhanced, breakable, colback=acento!4, colframe=acento,
      boxrule=0pt, leftrule=3pt, arc=0pt, title=Enunciado de la tarea,
      coltitle=acento, colbacktitle=acento!4, fonttitle=\bfseries]
    \enunciado
  \end{tcolorbox}
  \vspace{1em}
}

% ---------- Contenido ----------
\section*{Resolución}

\subsection*{Paso 1. Plantear la ecuación}
Cinco barras $x$ y un 12 a la izquierda; 67 a la derecha.
\captura{1.png}{$5x + 12 = 67$}

\subsection*{Paso 2. Añadir $-12$ a ambos lados}
Para eliminar el 12 se suma su opuesto en los dos platillos.
\captura{2.png}{$5x + 12 + (-12) = 67 + (-12)$}

\subsection*{Paso 3. Simplificar}
$12 + (-12) = 0$ y $67 + (-12) = 55$. El 0 se retira.
\captura{3.png}{$5x = 55$}

\subsection*{Paso 4.1. Descomponer 55}
Se descompone 55 en factores: $55 = 5 \cdot 11$.
\captura{4.1.png}{$55 = 5 \cdot 11$}

\subsection*{Paso 4.2. Separar los factores}
El 5 coincide con el número de barras $x$, así que a cada $x$ le corresponde un 11.
\captura{4.2.png}{$5x = 5 \cdot 11$}

\subsection*{Paso 5. Cambiar 55 por cinco fichas de 11}
\captura{5.png}{$x + x + x + x + x = 11 + 11 + 11 + 11 + 11$}

\subsection*{Paso 6. Quitar cantidades distintas}
Si se quitan cuatro $x$ pero solo tres 11, la balanza se desequilibra.
\captura{6.png}{$x < 22$}

\subsection*{Paso 7. Quitar lo mismo en ambos lados}
Se quitan cuatro $x$ y cuatro 11: la balanza vuelve a equilibrarse.
\captura{7.png}{$x = 11$}

\subsection*{Comprobación}
$5 \cdot 11 + 12 = 55 + 12 = 67$ \checkmark

\end{document}
